
% =====================================================================
% Results — consolidated experimental results for the two worked
% examples. Every figure is taken directly from the measured matrices:
%   - Curve25519: ../../curve25519/RESULTS.md
%   - Keccak:     ../bench/KECCAK_RESULTS.md
%
% Both campaigns share one measurement discipline: the Goldmont core
% (Intel Celeron N3350) is pinned to its ~1.09 GHz base frequency so the
% cycle counter and core clock advance together, every contender is timed
% back-to-back within a single process so all observe one clock, and each
% contender is swept over the 24 combinations of compiler and -O level
% with the best result retained, in the manner of the SUPERCOP autotuner.
%
% To preview, add  
% =====================================================================
% Results — consolidated experimental results for the two worked
% examples. Every figure is taken directly from the measured matrices:
%   - Curve25519: ../../curve25519/RESULTS.md
%   - Keccak:     ../bench/KECCAK_RESULTS.md
%
% Both campaigns share one measurement discipline: the Goldmont core
% (Intel Celeron N3350) is pinned to its ~1.09 GHz base frequency so the
% cycle counter and core clock advance together, every contender is timed
% back-to-back within a single process so all observe one clock, and each
% contender is swept over the 24 combinations of compiler and -O level
% with the best result retained, in the manner of the SUPERCOP autotuner.
%
% To preview, add  
% =====================================================================
% Results — consolidated experimental results for the two worked
% examples. Every figure is taken directly from the measured matrices:
%   - Curve25519: ../../curve25519/RESULTS.md
%   - Keccak:     ../bench/KECCAK_RESULTS.md
%
% Both campaigns share one measurement discipline: the Goldmont core
% (Intel Celeron N3350) is pinned to its ~1.09 GHz base frequency so the
% cycle counter and core clock advance together, every contender is timed
% back-to-back within a single process so all observe one clock, and each
% contender is swept over the 24 combinations of compiler and -O level
% with the best result retained, in the manner of the SUPERCOP autotuner.
%
% To preview, add  
% =====================================================================
% Results — consolidated experimental results for the two worked
% examples. Every figure is taken directly from the measured matrices:
%   - Curve25519: ../../curve25519/RESULTS.md
%   - Keccak:     ../bench/KECCAK_RESULTS.md
%
% Both campaigns share one measurement discipline: the Goldmont core
% (Intel Celeron N3350) is pinned to its ~1.09 GHz base frequency so the
% cycle counter and core clock advance together, every contender is timed
% back-to-back within a single process so all observe one clock, and each
% contender is swept over the 24 combinations of compiler and -O level
% with the best result retained, in the manner of the SUPERCOP autotuner.
%
% To preview, add  \input{section_results.tex}  to preview.tex.
% =====================================================================

\section{Results}
\label{sec:results}

This section gathers the measurements for the two worked examples under one
common discipline. All timings are taken on the Goldmont core pinned to its
$\sim\!1.09$\,GHz base frequency, so that the cycle counter advances with the
core clock; every contender is timed back to back within a single process so
that all observe the same clock; and each contender is swept over the
twenty-four combinations of compiler and optimisation level, the best result
being retained for each, in the manner of the SUPERCOP autotuner.

% ---------------------------------------------------------------------
\subsection{Curve25519}
\label{sec:res-c25519}

Timed in isolation in a tight loop, the two microcode field operations are the
fastest of their kind we measured: the squaring completes in roughly $37$ cycles,
about $26\%$ faster than the fiat-crypto reference, and the multiplication in
roughly $58$ cycles, about $3\%$ faster. The asymmetry follows from the kernels
themselves, since squaring issues only fifteen partial products to
multiplication's twenty-five and so leaves more slack into which the carry work
can be packed. End to end, the inlined microcode ladder completes a scalar
multiplication in about $299{,}967$ median cycles and is the fastest of every
$5\times51$ contender in the sweep, as \Cref{tab:res-c25519-e2e} shows, ahead of
\texttt{donna\_c64}, fiat-crypto, hand-written \texttt{\_\_uint128\_t} C, the
CryptOpt superoptimiser, and the Bernstein--Schwabe \texttt{amd64-51} hand
assembly.

\begin{table}[t]
\centering
\caption{End-to-end X25519: best configuration per contender across the
twenty-four-config sweep, in median cycles per scalar multiplication. Smaller is
better. The inlined microcode ladder is the fastest $5\times51$ implementation but
is overtaken by the hand-written $4\times64$ saturated assembly, and the microcode
port to that same $4\times64$ representation is the slowest entry by a wide margin.}
\label{tab:res-c25519-e2e}
\small
\begin{tabular}{@{}llr@{}}
\toprule
contender & representation / backend & median (cyc) \\
\midrule
amd64-64 / asm       & $4\times64$ saturated, hand asm                     & 271\,326 \\
ours / ucode-inline  & $5\times51$, microcode (inlined ladder)            & 299\,967 \\
amd64-51 / ucode     & $5\times51$, microcode (B--S framework)            & 314\,416 \\
donna\_c64           & $5\times51$, portable C                            & 335\,744 \\
ours / fiat          & $5\times51$, fiat-crypto C                         & 354\,851 \\
amd64-51 / asm       & $5\times51$, hand asm                              & 356\,038 \\
ours / hand-C        & $5\times51$, hand-written \texttt{\_\_uint128\_t} C & 379\,433 \\
ours / cryptopt      & $5\times51$, CryptOpt asm                          & 380\,361 \\
amd64-64 / ucode     & $4\times64$ saturated, microcode                   & 511\,017 \\
\bottomrule
\end{tabular}
\end{table}

\Cref{tab:res-c25519-ratio} aggregates the whole sweep as the geometric mean of
the per-configuration ratio against the inlined ladder. Across the twenty-four
configurations it is a geomean $1.25\times$ faster than fiat-crypto, $1.28\times$
faster than CryptOpt, $1.30\times$ faster than hand-written C, and $1.31\times$
faster than \texttt{donna\_c64}, and the field-operation advantage is not an
artefact of our own ladder: dropped unchanged into the \texttt{amd64-51}
framework it makes that hybrid a geomean $1.08\times$ faster than the same
framework running its authors' own assembly field operations. The limit is
equally clear. The fastest X25519 on this core is the saturated $4\times64$
\texttt{amd64-64} assembly at $271{,}326$ cycles, roughly $11\%$ ahead of the
inlined ladder (geomean $0.89\times$), and porting the field operations into that
same $4\times64$ representation produces the slowest entry in the table at
$511{,}017$ cycles --- a geomean $1.77\times$ slower than our $5\times51$
microcode --- because the long, full-width carry chains of a saturated
multiplication must be emulated by serial \texttt{SETCC} operations under the
frozen single-flag model. The representation that suits the native core best is
thus precisely the one that suits the microcode layer worst.

\begin{table}[t]
\centering
\caption{Geometric-mean speedup of the inlined microcode ladder
(\texttt{ours/ucode-inline}) over each contender across the twenty-four-config
sweep, computed as the geomean of the per-config ratio (contender median $\div$
ucode-inline median). A value $>1$ means the microcode ladder is faster.}
\label{tab:res-c25519-ratio}
\small
\begin{tabular}{@{}lr@{}}
\toprule
contender & geomean speedup of ucode-inline \\
\midrule
amd64-64 / ucode  & $1.77\times$ \\
donna\_c64        & $1.31\times$ \\
ours / hand-C     & $1.30\times$ \\
ours / cryptopt   & $1.28\times$ \\
ours / fiat       & $1.25\times$ \\
amd64-51 / asm    & $1.17\times$ \\
amd64-51 / ucode  & $1.08\times$ \\
amd64-64 / asm    & $0.89\times$ \\
\bottomrule
\end{tabular}
\end{table}

% ---------------------------------------------------------------------
\subsection{Keccak-$f[1600]$}
\label{sec:res-keccak}

Under the same base-pinned, same-process measurement, the looped permutation
completes in $1910$ cycles per permutation at its best configuration (gcc-12
\texttt{-Os}) and is the fastest of all ten contenders, as
\Cref{tab:res-keccak-best} reports. The nearest scalar baselines are the SUPERCOP
autotuner's pick \texttt{opt64lcu24} at $2047$\,cyc/perm and the hand-tuned
\texttt{x86\_64\_asm} at $2062$, giving a ratio of $0.933$ --- an improvement of
approximately $7\%$. The ratio, rather than the absolute count, is the robust
quantity: \texttt{rdtsc} ticks at the $\sim\!1.1$\,GHz reference rate while the
core bursts to $\sim\!2.4$\,GHz, so the same harness run hot reports
proportionally fewer ticks but an identical ratio. The two \texttt{shld}-based
variants are roughly five times slower because that rotate idiom is heavily
penalised on this microarchitecture, and they are reported only for completeness.

\begin{table}[t]
\centering
\caption{Keccak-$f[1600]$: best result per contender across the
twenty-four-config sweep, in cycles per permutation, base-pinned and measured
same-process. Smaller is better.}
\label{tab:res-keccak-best}
\small
\begin{tabular}{@{}lrrl@{}}
\toprule
contender & min (cyc/perm) & median & best config \\
\midrule
\textbf{microcode (looped)} & \textbf{1910} & 1915  & gcc-12 -Os \\
keccak / opt64lcu24         & 2047          & 2056  & gcc-11 -Os \\
keccak / x86\_64\_asm       & 2062          & 2069  & gcc-11 -O2 \\
keccak / opt64lcu6          & 2204          & 2207  & clang-14 -O \\
keccak / opt64u6            & 2400          & 2402  & gcc-11 -O2 \\
keccak / simple             & 2493          & 2495  & clang-17 -Os \\
keccak / sseu2              & 3055          & 3064  & clang-18 -Os \\
keccak / mmxu1              & 3867          & 3873  & gcc-12 -O \\
keccak / opt64lcu24shld     & 10053         & 10061 & clang-17 -Os \\
keccak / x86\_64\_shld      & 10093         & 10101 & gcc-13 -Os \\
\bottomrule
\end{tabular}
\end{table}

Beyond winning at its best configuration, the microcode wins at every one of the
twenty-four configurations and is among the most stable contenders in the sweep,
as \Cref{tab:res-keccak-stability} shows: its cost varies by only $5$ cycles
across all compilers and optimisation levels ($1910$--$1915$), whereas the fastest
scalar baseline, \texttt{opt64lcu24}, swings by $211$ cycles ($2047$--$2258$) and
attains its minimum only when a particular compiler happens to schedule it well.
This stability is a direct consequence of the design: a single looped round body
executed entirely within microcode is insensitive to the host compiler, which
touches only the thin native prologue that fires the patch. The margin itself is
modest by design and bounded by latency. The permutation is a strict twenty-four
round dependency chain with a backward branch per round and no overlap between
firings, so it pays the true per-triad latency of roughly $1.37$ cycles over its
$1396$ dynamically executed triads, and the 128-triad cap precludes the unrolling
that would expose independent work. The entire advantage is therefore structural
--- full register residency across all $24$ rounds, the borrowed $32$nd register
that keeps the $\theta$ mixing lanes resident, and the move-free $\chi$ stage ---
rather than a higher peak throughput.

\begin{table}[t]
\centering
\caption{Stability across the twenty-four-config sweep for the competitive
64-bit contenders: the cost at the best configuration against the worst. The
microcode is both the fastest and among the most stable, varying by only $5$
cycles, whereas the fastest scalar baseline, \texttt{opt64lcu24}, swings by over
$200$ cycles depending on the compiler.}
\label{tab:res-keccak-stability}
\small
\begin{tabular}{@{}lrrr@{}}
\toprule
contender & best (cyc/perm) & worst & spread \\
\midrule
microcode (looped)    & 1910 & 1915 & 5   \\
keccak / x86\_64\_asm & 2062 & 2066 & 4   \\
keccak / opt64lcu24   & 2047 & 2258 & 211 \\
keccak / opt64lcu6    & 2204 & 2209 & 5   \\
keccak / opt64u6      & 2400 & 2403 & 3   \\
keccak / simple       & 2493 & 2513 & 20  \\
keccak / sseu2        & 3055 & 3137 & 82  \\
keccak / mmxu1        & 3867 & 3878 & 11  \\
\bottomrule
\end{tabular}
\end{table}
  to preview.tex.
% =====================================================================

\section{Results}
\label{sec:results}

This section gathers the measurements for the two worked examples under one
common discipline. All timings are taken on the Goldmont core pinned to its
$\sim\!1.09$\,GHz base frequency, so that the cycle counter advances with the
core clock; every contender is timed back to back within a single process so
that all observe the same clock; and each contender is swept over the
twenty-four combinations of compiler and optimisation level, the best result
being retained for each, in the manner of the SUPERCOP autotuner.

% ---------------------------------------------------------------------
\subsection{Curve25519}
\label{sec:res-c25519}

Timed in isolation in a tight loop, the two microcode field operations are the
fastest of their kind we measured: the squaring completes in roughly $37$ cycles,
about $26\%$ faster than the fiat-crypto reference, and the multiplication in
roughly $58$ cycles, about $3\%$ faster. The asymmetry follows from the kernels
themselves, since squaring issues only fifteen partial products to
multiplication's twenty-five and so leaves more slack into which the carry work
can be packed. End to end, the inlined microcode ladder completes a scalar
multiplication in about $299{,}967$ median cycles and is the fastest of every
$5\times51$ contender in the sweep, as \Cref{tab:res-c25519-e2e} shows, ahead of
\texttt{donna\_c64}, fiat-crypto, hand-written \texttt{\_\_uint128\_t} C, the
CryptOpt superoptimiser, and the Bernstein--Schwabe \texttt{amd64-51} hand
assembly.

\begin{table}[t]
\centering
\caption{End-to-end X25519: best configuration per contender across the
twenty-four-config sweep, in median cycles per scalar multiplication. Smaller is
better. The inlined microcode ladder is the fastest $5\times51$ implementation but
is overtaken by the hand-written $4\times64$ saturated assembly, and the microcode
port to that same $4\times64$ representation is the slowest entry by a wide margin.}
\label{tab:res-c25519-e2e}
\small
\begin{tabular}{@{}llr@{}}
\toprule
contender & representation / backend & median (cyc) \\
\midrule
amd64-64 / asm       & $4\times64$ saturated, hand asm                     & 271\,326 \\
ours / ucode-inline  & $5\times51$, microcode (inlined ladder)            & 299\,967 \\
amd64-51 / ucode     & $5\times51$, microcode (B--S framework)            & 314\,416 \\
donna\_c64           & $5\times51$, portable C                            & 335\,744 \\
ours / fiat          & $5\times51$, fiat-crypto C                         & 354\,851 \\
amd64-51 / asm       & $5\times51$, hand asm                              & 356\,038 \\
ours / hand-C        & $5\times51$, hand-written \texttt{\_\_uint128\_t} C & 379\,433 \\
ours / cryptopt      & $5\times51$, CryptOpt asm                          & 380\,361 \\
amd64-64 / ucode     & $4\times64$ saturated, microcode                   & 511\,017 \\
\bottomrule
\end{tabular}
\end{table}

\Cref{tab:res-c25519-ratio} aggregates the whole sweep as the geometric mean of
the per-configuration ratio against the inlined ladder. Across the twenty-four
configurations it is a geomean $1.25\times$ faster than fiat-crypto, $1.28\times$
faster than CryptOpt, $1.30\times$ faster than hand-written C, and $1.31\times$
faster than \texttt{donna\_c64}, and the field-operation advantage is not an
artefact of our own ladder: dropped unchanged into the \texttt{amd64-51}
framework it makes that hybrid a geomean $1.08\times$ faster than the same
framework running its authors' own assembly field operations. The limit is
equally clear. The fastest X25519 on this core is the saturated $4\times64$
\texttt{amd64-64} assembly at $271{,}326$ cycles, roughly $11\%$ ahead of the
inlined ladder (geomean $0.89\times$), and porting the field operations into that
same $4\times64$ representation produces the slowest entry in the table at
$511{,}017$ cycles --- a geomean $1.77\times$ slower than our $5\times51$
microcode --- because the long, full-width carry chains of a saturated
multiplication must be emulated by serial \texttt{SETCC} operations under the
frozen single-flag model. The representation that suits the native core best is
thus precisely the one that suits the microcode layer worst.

\begin{table}[t]
\centering
\caption{Geometric-mean speedup of the inlined microcode ladder
(\texttt{ours/ucode-inline}) over each contender across the twenty-four-config
sweep, computed as the geomean of the per-config ratio (contender median $\div$
ucode-inline median). A value $>1$ means the microcode ladder is faster.}
\label{tab:res-c25519-ratio}
\small
\begin{tabular}{@{}lr@{}}
\toprule
contender & geomean speedup of ucode-inline \\
\midrule
amd64-64 / ucode  & $1.77\times$ \\
donna\_c64        & $1.31\times$ \\
ours / hand-C     & $1.30\times$ \\
ours / cryptopt   & $1.28\times$ \\
ours / fiat       & $1.25\times$ \\
amd64-51 / asm    & $1.17\times$ \\
amd64-51 / ucode  & $1.08\times$ \\
amd64-64 / asm    & $0.89\times$ \\
\bottomrule
\end{tabular}
\end{table}

% ---------------------------------------------------------------------
\subsection{Keccak-$f[1600]$}
\label{sec:res-keccak}

Under the same base-pinned, same-process measurement, the looped permutation
completes in $1910$ cycles per permutation at its best configuration (gcc-12
\texttt{-Os}) and is the fastest of all ten contenders, as
\Cref{tab:res-keccak-best} reports. The nearest scalar baselines are the SUPERCOP
autotuner's pick \texttt{opt64lcu24} at $2047$\,cyc/perm and the hand-tuned
\texttt{x86\_64\_asm} at $2062$, giving a ratio of $0.933$ --- an improvement of
approximately $7\%$. The ratio, rather than the absolute count, is the robust
quantity: \texttt{rdtsc} ticks at the $\sim\!1.1$\,GHz reference rate while the
core bursts to $\sim\!2.4$\,GHz, so the same harness run hot reports
proportionally fewer ticks but an identical ratio. The two \texttt{shld}-based
variants are roughly five times slower because that rotate idiom is heavily
penalised on this microarchitecture, and they are reported only for completeness.

\begin{table}[t]
\centering
\caption{Keccak-$f[1600]$: best result per contender across the
twenty-four-config sweep, in cycles per permutation, base-pinned and measured
same-process. Smaller is better.}
\label{tab:res-keccak-best}
\small
\begin{tabular}{@{}lrrl@{}}
\toprule
contender & min (cyc/perm) & median & best config \\
\midrule
\textbf{microcode (looped)} & \textbf{1910} & 1915  & gcc-12 -Os \\
keccak / opt64lcu24         & 2047          & 2056  & gcc-11 -Os \\
keccak / x86\_64\_asm       & 2062          & 2069  & gcc-11 -O2 \\
keccak / opt64lcu6          & 2204          & 2207  & clang-14 -O \\
keccak / opt64u6            & 2400          & 2402  & gcc-11 -O2 \\
keccak / simple             & 2493          & 2495  & clang-17 -Os \\
keccak / sseu2              & 3055          & 3064  & clang-18 -Os \\
keccak / mmxu1              & 3867          & 3873  & gcc-12 -O \\
keccak / opt64lcu24shld     & 10053         & 10061 & clang-17 -Os \\
keccak / x86\_64\_shld      & 10093         & 10101 & gcc-13 -Os \\
\bottomrule
\end{tabular}
\end{table}

Beyond winning at its best configuration, the microcode wins at every one of the
twenty-four configurations and is among the most stable contenders in the sweep,
as \Cref{tab:res-keccak-stability} shows: its cost varies by only $5$ cycles
across all compilers and optimisation levels ($1910$--$1915$), whereas the fastest
scalar baseline, \texttt{opt64lcu24}, swings by $211$ cycles ($2047$--$2258$) and
attains its minimum only when a particular compiler happens to schedule it well.
This stability is a direct consequence of the design: a single looped round body
executed entirely within microcode is insensitive to the host compiler, which
touches only the thin native prologue that fires the patch. The margin itself is
modest by design and bounded by latency. The permutation is a strict twenty-four
round dependency chain with a backward branch per round and no overlap between
firings, so it pays the true per-triad latency of roughly $1.37$ cycles over its
$1396$ dynamically executed triads, and the 128-triad cap precludes the unrolling
that would expose independent work. The entire advantage is therefore structural
--- full register residency across all $24$ rounds, the borrowed $32$nd register
that keeps the $\theta$ mixing lanes resident, and the move-free $\chi$ stage ---
rather than a higher peak throughput.

\begin{table}[t]
\centering
\caption{Stability across the twenty-four-config sweep for the competitive
64-bit contenders: the cost at the best configuration against the worst. The
microcode is both the fastest and among the most stable, varying by only $5$
cycles, whereas the fastest scalar baseline, \texttt{opt64lcu24}, swings by over
$200$ cycles depending on the compiler.}
\label{tab:res-keccak-stability}
\small
\begin{tabular}{@{}lrrr@{}}
\toprule
contender & best (cyc/perm) & worst & spread \\
\midrule
microcode (looped)    & 1910 & 1915 & 5   \\
keccak / x86\_64\_asm & 2062 & 2066 & 4   \\
keccak / opt64lcu24   & 2047 & 2258 & 211 \\
keccak / opt64lcu6    & 2204 & 2209 & 5   \\
keccak / opt64u6      & 2400 & 2403 & 3   \\
keccak / simple       & 2493 & 2513 & 20  \\
keccak / sseu2        & 3055 & 3137 & 82  \\
keccak / mmxu1        & 3867 & 3878 & 11  \\
\bottomrule
\end{tabular}
\end{table}
  to preview.tex.
% =====================================================================

\section{Results}
\label{sec:results}

This section gathers the measurements for the two worked examples under one
common discipline. All timings are taken on the Goldmont core pinned to its
$\sim\!1.09$\,GHz base frequency, so that the cycle counter advances with the
core clock; every contender is timed back to back within a single process so
that all observe the same clock; and each contender is swept over the
twenty-four combinations of compiler and optimisation level, the best result
being retained for each, in the manner of the SUPERCOP autotuner.

% ---------------------------------------------------------------------
\subsection{Curve25519}
\label{sec:res-c25519}

Timed in isolation in a tight loop, the two microcode field operations are the
fastest of their kind we measured: the squaring completes in roughly $37$ cycles,
about $26\%$ faster than the fiat-crypto reference, and the multiplication in
roughly $58$ cycles, about $3\%$ faster. The asymmetry follows from the kernels
themselves, since squaring issues only fifteen partial products to
multiplication's twenty-five and so leaves more slack into which the carry work
can be packed. End to end, the inlined microcode ladder completes a scalar
multiplication in about $299{,}967$ median cycles and is the fastest of every
$5\times51$ contender in the sweep, as \Cref{tab:res-c25519-e2e} shows, ahead of
\texttt{donna\_c64}, fiat-crypto, hand-written \texttt{\_\_uint128\_t} C, the
CryptOpt superoptimiser, and the Bernstein--Schwabe \texttt{amd64-51} hand
assembly.

\begin{table}[t]
\centering
\caption{End-to-end X25519: best configuration per contender across the
twenty-four-config sweep, in median cycles per scalar multiplication. Smaller is
better. The inlined microcode ladder is the fastest $5\times51$ implementation but
is overtaken by the hand-written $4\times64$ saturated assembly, and the microcode
port to that same $4\times64$ representation is the slowest entry by a wide margin.}
\label{tab:res-c25519-e2e}
\small
\begin{tabular}{@{}llr@{}}
\toprule
contender & representation / backend & median (cyc) \\
\midrule
amd64-64 / asm       & $4\times64$ saturated, hand asm                     & 271\,326 \\
ours / ucode-inline  & $5\times51$, microcode (inlined ladder)            & 299\,967 \\
amd64-51 / ucode     & $5\times51$, microcode (B--S framework)            & 314\,416 \\
donna\_c64           & $5\times51$, portable C                            & 335\,744 \\
ours / fiat          & $5\times51$, fiat-crypto C                         & 354\,851 \\
amd64-51 / asm       & $5\times51$, hand asm                              & 356\,038 \\
ours / hand-C        & $5\times51$, hand-written \texttt{\_\_uint128\_t} C & 379\,433 \\
ours / cryptopt      & $5\times51$, CryptOpt asm                          & 380\,361 \\
amd64-64 / ucode     & $4\times64$ saturated, microcode                   & 511\,017 \\
\bottomrule
\end{tabular}
\end{table}

\Cref{tab:res-c25519-ratio} aggregates the whole sweep as the geometric mean of
the per-configuration ratio against the inlined ladder. Across the twenty-four
configurations it is a geomean $1.25\times$ faster than fiat-crypto, $1.28\times$
faster than CryptOpt, $1.30\times$ faster than hand-written C, and $1.31\times$
faster than \texttt{donna\_c64}, and the field-operation advantage is not an
artefact of our own ladder: dropped unchanged into the \texttt{amd64-51}
framework it makes that hybrid a geomean $1.08\times$ faster than the same
framework running its authors' own assembly field operations. The limit is
equally clear. The fastest X25519 on this core is the saturated $4\times64$
\texttt{amd64-64} assembly at $271{,}326$ cycles, roughly $11\%$ ahead of the
inlined ladder (geomean $0.89\times$), and porting the field operations into that
same $4\times64$ representation produces the slowest entry in the table at
$511{,}017$ cycles --- a geomean $1.77\times$ slower than our $5\times51$
microcode --- because the long, full-width carry chains of a saturated
multiplication must be emulated by serial \texttt{SETCC} operations under the
frozen single-flag model. The representation that suits the native core best is
thus precisely the one that suits the microcode layer worst.

\begin{table}[t]
\centering
\caption{Geometric-mean speedup of the inlined microcode ladder
(\texttt{ours/ucode-inline}) over each contender across the twenty-four-config
sweep, computed as the geomean of the per-config ratio (contender median $\div$
ucode-inline median). A value $>1$ means the microcode ladder is faster.}
\label{tab:res-c25519-ratio}
\small
\begin{tabular}{@{}lr@{}}
\toprule
contender & geomean speedup of ucode-inline \\
\midrule
amd64-64 / ucode  & $1.77\times$ \\
donna\_c64        & $1.31\times$ \\
ours / hand-C     & $1.30\times$ \\
ours / cryptopt   & $1.28\times$ \\
ours / fiat       & $1.25\times$ \\
amd64-51 / asm    & $1.17\times$ \\
amd64-51 / ucode  & $1.08\times$ \\
amd64-64 / asm    & $0.89\times$ \\
\bottomrule
\end{tabular}
\end{table}

% ---------------------------------------------------------------------
\subsection{Keccak-$f[1600]$}
\label{sec:res-keccak}

Under the same base-pinned, same-process measurement, the looped permutation
completes in $1910$ cycles per permutation at its best configuration (gcc-12
\texttt{-Os}) and is the fastest of all ten contenders, as
\Cref{tab:res-keccak-best} reports. The nearest scalar baselines are the SUPERCOP
autotuner's pick \texttt{opt64lcu24} at $2047$\,cyc/perm and the hand-tuned
\texttt{x86\_64\_asm} at $2062$, giving a ratio of $0.933$ --- an improvement of
approximately $7\%$. The ratio, rather than the absolute count, is the robust
quantity: \texttt{rdtsc} ticks at the $\sim\!1.1$\,GHz reference rate while the
core bursts to $\sim\!2.4$\,GHz, so the same harness run hot reports
proportionally fewer ticks but an identical ratio. The two \texttt{shld}-based
variants are roughly five times slower because that rotate idiom is heavily
penalised on this microarchitecture, and they are reported only for completeness.

\begin{table}[t]
\centering
\caption{Keccak-$f[1600]$: best result per contender across the
twenty-four-config sweep, in cycles per permutation, base-pinned and measured
same-process. Smaller is better.}
\label{tab:res-keccak-best}
\small
\begin{tabular}{@{}lrrl@{}}
\toprule
contender & min (cyc/perm) & median & best config \\
\midrule
\textbf{microcode (looped)} & \textbf{1910} & 1915  & gcc-12 -Os \\
keccak / opt64lcu24         & 2047          & 2056  & gcc-11 -Os \\
keccak / x86\_64\_asm       & 2062          & 2069  & gcc-11 -O2 \\
keccak / opt64lcu6          & 2204          & 2207  & clang-14 -O \\
keccak / opt64u6            & 2400          & 2402  & gcc-11 -O2 \\
keccak / simple             & 2493          & 2495  & clang-17 -Os \\
keccak / sseu2              & 3055          & 3064  & clang-18 -Os \\
keccak / mmxu1              & 3867          & 3873  & gcc-12 -O \\
keccak / opt64lcu24shld     & 10053         & 10061 & clang-17 -Os \\
keccak / x86\_64\_shld      & 10093         & 10101 & gcc-13 -Os \\
\bottomrule
\end{tabular}
\end{table}

Beyond winning at its best configuration, the microcode wins at every one of the
twenty-four configurations and is among the most stable contenders in the sweep,
as \Cref{tab:res-keccak-stability} shows: its cost varies by only $5$ cycles
across all compilers and optimisation levels ($1910$--$1915$), whereas the fastest
scalar baseline, \texttt{opt64lcu24}, swings by $211$ cycles ($2047$--$2258$) and
attains its minimum only when a particular compiler happens to schedule it well.
This stability is a direct consequence of the design: a single looped round body
executed entirely within microcode is insensitive to the host compiler, which
touches only the thin native prologue that fires the patch. The margin itself is
modest by design and bounded by latency. The permutation is a strict twenty-four
round dependency chain with a backward branch per round and no overlap between
firings, so it pays the true per-triad latency of roughly $1.37$ cycles over its
$1396$ dynamically executed triads, and the 128-triad cap precludes the unrolling
that would expose independent work. The entire advantage is therefore structural
--- full register residency across all $24$ rounds, the borrowed $32$nd register
that keeps the $\theta$ mixing lanes resident, and the move-free $\chi$ stage ---
rather than a higher peak throughput.

\begin{table}[t]
\centering
\caption{Stability across the twenty-four-config sweep for the competitive
64-bit contenders: the cost at the best configuration against the worst. The
microcode is both the fastest and among the most stable, varying by only $5$
cycles, whereas the fastest scalar baseline, \texttt{opt64lcu24}, swings by over
$200$ cycles depending on the compiler.}
\label{tab:res-keccak-stability}
\small
\begin{tabular}{@{}lrrr@{}}
\toprule
contender & best (cyc/perm) & worst & spread \\
\midrule
microcode (looped)    & 1910 & 1915 & 5   \\
keccak / x86\_64\_asm & 2062 & 2066 & 4   \\
keccak / opt64lcu24   & 2047 & 2258 & 211 \\
keccak / opt64lcu6    & 2204 & 2209 & 5   \\
keccak / opt64u6      & 2400 & 2403 & 3   \\
keccak / simple       & 2493 & 2513 & 20  \\
keccak / sseu2        & 3055 & 3137 & 82  \\
keccak / mmxu1        & 3867 & 3878 & 11  \\
\bottomrule
\end{tabular}
\end{table}
  to preview.tex.
% =====================================================================

\section{Results}
\label{sec:results}

This section gathers the measurements for the two worked examples under one
common discipline. All timings are taken on the Goldmont core pinned to its
$\sim\!1.09$\,GHz base frequency, so that the cycle counter advances with the
core clock; every contender is timed back to back within a single process so
that all observe the same clock; and each contender is swept over the
twenty-four combinations of compiler and optimisation level, the best result
being retained for each, in the manner of the SUPERCOP autotuner.

% ---------------------------------------------------------------------
\subsection{Curve25519}
\label{sec:res-c25519}

Timed in isolation in a tight loop, the two microcode field operations are the
fastest of their kind we measured: the squaring completes in roughly $37$ cycles,
about $26\%$ faster than the fiat-crypto reference, and the multiplication in
roughly $58$ cycles, about $3\%$ faster. The asymmetry follows from the kernels
themselves, since squaring issues only fifteen partial products to
multiplication's twenty-five and so leaves more slack into which the carry work
can be packed. End to end, the inlined microcode ladder completes a scalar
multiplication in about $299{,}967$ median cycles and is the fastest of every
$5\times51$ contender in the sweep, as \Cref{tab:res-c25519-e2e} shows, ahead of
\texttt{donna\_c64}, fiat-crypto, hand-written \texttt{\_\_uint128\_t} C, the
CryptOpt superoptimiser, and the Bernstein--Schwabe \texttt{amd64-51} hand
assembly.

\begin{table}[t]
\centering
\caption{End-to-end X25519: best configuration per contender across the
twenty-four-config sweep, in median cycles per scalar multiplication. Smaller is
better. The inlined microcode ladder is the fastest $5\times51$ implementation but
is overtaken by the hand-written $4\times64$ saturated assembly, and the microcode
port to that same $4\times64$ representation is the slowest entry by a wide margin.}
\label{tab:res-c25519-e2e}
\small
\begin{tabular}{@{}llr@{}}
\toprule
contender & representation / backend & median (cyc) \\
\midrule
amd64-64 / asm       & $4\times64$ saturated, hand asm                     & 271\,326 \\
ours / ucode-inline  & $5\times51$, microcode (inlined ladder)            & 299\,967 \\
amd64-51 / ucode     & $5\times51$, microcode (B--S framework)            & 314\,416 \\
donna\_c64           & $5\times51$, portable C                            & 335\,744 \\
ours / fiat          & $5\times51$, fiat-crypto C                         & 354\,851 \\
amd64-51 / asm       & $5\times51$, hand asm                              & 356\,038 \\
ours / hand-C        & $5\times51$, hand-written \texttt{\_\_uint128\_t} C & 379\,433 \\
ours / cryptopt      & $5\times51$, CryptOpt asm                          & 380\,361 \\
amd64-64 / ucode     & $4\times64$ saturated, microcode                   & 511\,017 \\
\bottomrule
\end{tabular}
\end{table}

\Cref{tab:res-c25519-ratio} aggregates the whole sweep as the geometric mean of
the per-configuration ratio against the inlined ladder. Across the twenty-four
configurations it is a geomean $1.25\times$ faster than fiat-crypto, $1.28\times$
faster than CryptOpt, $1.30\times$ faster than hand-written C, and $1.31\times$
faster than \texttt{donna\_c64}, and the field-operation advantage is not an
artefact of our own ladder: dropped unchanged into the \texttt{amd64-51}
framework it makes that hybrid a geomean $1.08\times$ faster than the same
framework running its authors' own assembly field operations. The limit is
equally clear. The fastest X25519 on this core is the saturated $4\times64$
\texttt{amd64-64} assembly at $271{,}326$ cycles, roughly $11\%$ ahead of the
inlined ladder (geomean $0.89\times$), and porting the field operations into that
same $4\times64$ representation produces the slowest entry in the table at
$511{,}017$ cycles --- a geomean $1.77\times$ slower than our $5\times51$
microcode --- because the long, full-width carry chains of a saturated
multiplication must be emulated by serial \texttt{SETCC} operations under the
frozen single-flag model. The representation that suits the native core best is
thus precisely the one that suits the microcode layer worst.

\begin{table}[t]
\centering
\caption{Geometric-mean speedup of the inlined microcode ladder
(\texttt{ours/ucode-inline}) over each contender across the twenty-four-config
sweep, computed as the geomean of the per-config ratio (contender median $\div$
ucode-inline median). A value $>1$ means the microcode ladder is faster.}
\label{tab:res-c25519-ratio}
\small
\begin{tabular}{@{}lr@{}}
\toprule
contender & geomean speedup of ucode-inline \\
\midrule
amd64-64 / ucode  & $1.77\times$ \\
donna\_c64        & $1.31\times$ \\
ours / hand-C     & $1.30\times$ \\
ours / cryptopt   & $1.28\times$ \\
ours / fiat       & $1.25\times$ \\
amd64-51 / asm    & $1.17\times$ \\
amd64-51 / ucode  & $1.08\times$ \\
amd64-64 / asm    & $0.89\times$ \\
\bottomrule
\end{tabular}
\end{table}

% ---------------------------------------------------------------------
\subsection{Keccak-$f[1600]$}
\label{sec:res-keccak}

Under the same base-pinned, same-process measurement, the looped permutation
completes in $1910$ cycles per permutation at its best configuration (gcc-12
\texttt{-Os}) and is the fastest of all ten contenders, as
\Cref{tab:res-keccak-best} reports. The nearest scalar baselines are the SUPERCOP
autotuner's pick \texttt{opt64lcu24} at $2047$\,cyc/perm and the hand-tuned
\texttt{x86\_64\_asm} at $2062$, giving a ratio of $0.933$ --- an improvement of
approximately $7\%$. The ratio, rather than the absolute count, is the robust
quantity: \texttt{rdtsc} ticks at the $\sim\!1.1$\,GHz reference rate while the
core bursts to $\sim\!2.4$\,GHz, so the same harness run hot reports
proportionally fewer ticks but an identical ratio. The two \texttt{shld}-based
variants are roughly five times slower because that rotate idiom is heavily
penalised on this microarchitecture, and they are reported only for completeness.

\begin{table}[t]
\centering
\caption{Keccak-$f[1600]$: best result per contender across the
twenty-four-config sweep, in cycles per permutation, base-pinned and measured
same-process. Smaller is better.}
\label{tab:res-keccak-best}
\small
\begin{tabular}{@{}lrrl@{}}
\toprule
contender & min (cyc/perm) & median & best config \\
\midrule
\textbf{microcode (looped)} & \textbf{1910} & 1915  & gcc-12 -Os \\
keccak / opt64lcu24         & 2047          & 2056  & gcc-11 -Os \\
keccak / x86\_64\_asm       & 2062          & 2069  & gcc-11 -O2 \\
keccak / opt64lcu6          & 2204          & 2207  & clang-14 -O \\
keccak / opt64u6            & 2400          & 2402  & gcc-11 -O2 \\
keccak / simple             & 2493          & 2495  & clang-17 -Os \\
keccak / sseu2              & 3055          & 3064  & clang-18 -Os \\
keccak / mmxu1              & 3867          & 3873  & gcc-12 -O \\
keccak / opt64lcu24shld     & 10053         & 10061 & clang-17 -Os \\
keccak / x86\_64\_shld      & 10093         & 10101 & gcc-13 -Os \\
\bottomrule
\end{tabular}
\end{table}

Beyond winning at its best configuration, the microcode wins at every one of the
twenty-four configurations and is among the most stable contenders in the sweep,
as \Cref{tab:res-keccak-stability} shows: its cost varies by only $5$ cycles
across all compilers and optimisation levels ($1910$--$1915$), whereas the fastest
scalar baseline, \texttt{opt64lcu24}, swings by $211$ cycles ($2047$--$2258$) and
attains its minimum only when a particular compiler happens to schedule it well.
This stability is a direct consequence of the design: a single looped round body
executed entirely within microcode is insensitive to the host compiler, which
touches only the thin native prologue that fires the patch. The margin itself is
modest by design and bounded by latency. The permutation is a strict twenty-four
round dependency chain with a backward branch per round and no overlap between
firings, so it pays the true per-triad latency of roughly $1.37$ cycles over its
$1396$ dynamically executed triads, and the 128-triad cap precludes the unrolling
that would expose independent work. The entire advantage is therefore structural
--- full register residency across all $24$ rounds, the borrowed $32$nd register
that keeps the $\theta$ mixing lanes resident, and the move-free $\chi$ stage ---
rather than a higher peak throughput.

\begin{table}[t]
\centering
\caption{Stability across the twenty-four-config sweep for the competitive
64-bit contenders: the cost at the best configuration against the worst. The
microcode is both the fastest and among the most stable, varying by only $5$
cycles, whereas the fastest scalar baseline, \texttt{opt64lcu24}, swings by over
$200$ cycles depending on the compiler.}
\label{tab:res-keccak-stability}
\small
\begin{tabular}{@{}lrrr@{}}
\toprule
contender & best (cyc/perm) & worst & spread \\
\midrule
microcode (looped)    & 1910 & 1915 & 5   \\
keccak / x86\_64\_asm & 2062 & 2066 & 4   \\
keccak / opt64lcu24   & 2047 & 2258 & 211 \\
keccak / opt64lcu6    & 2204 & 2209 & 5   \\
keccak / opt64u6      & 2400 & 2403 & 3   \\
keccak / simple       & 2493 & 2513 & 20  \\
keccak / sseu2        & 3055 & 3137 & 82  \\
keccak / mmxu1        & 3867 & 3878 & 11  \\
\bottomrule
\end{tabular}
\end{table}
